\scanpage{5}
encore une fois que c'est le \uline{germe}
seulement des $\varphi_\xi$ qui est défini
(\unsure{avec variation à} ss groupe d'indice
fini). [Si $X$ s'immerge injectivement
dans une extension d'une VA par un
tore, alors $\xi$ est connu quand on
connaît ce germe
\[
(11)\qquad\varphi_\xi\in\underset{\substack{\Gamma'\subset\Gamma\\\text{ouvert}}}{\lim\limits_{\longrightarrow}}
\operatorname{Hom}(\Gamma',\mathcal E_{\deleted{\unsure{\overline X}}\inserted{X}})
/\operatorname{Int}(\pi_1(X),\mathcal E_X)
\]
[et même l'image de $\varphi_\xi$ dans un $H^1$ abélien
convenable\unsure{\ldots}] ]

\textcircled{4}
\[
(12)\qquad\text{Soit }\boxed{D=\sum_{i\in I}D_i}\quad\text{un diviseur \inserted{effectif} sur }X,
\]
($I$ l'ens de ses comp. \unsure{irré d|irred}) sans comp. mult.,
supposons que $X$ soit régulier \inserted{(i.e. lisse)} en les
pts gén. de chaque $D_i$. Soit \deleted{$\unsure{X|X'}$}
\[
(13)\qquad X'=X-\operatorname{supp}D
\]
et considérons le sous-groupe $\mathcal G_{X';X}$ de
$\mathcal G_X$ qui invarie $D$. On le notera \inserted{encore} (par
abus de notation) $\mathcal G_{X'}$, \inserted{ou plus précisément $\mathcal G_{X'}(I)$.} Alors on a
\mspage{4}
